\chapter{SciPy による数値計算とフィット}
\label{chap:scipy}

SciPy\index{scipy@SciPy|textbf} は、NumPy 配列を入力として統計、補間、信号処理、最適化、数値積分を行うライブラリ群である。本章では、分布モデル、パラメータ推定、曲線フィット、補間、信号処理と、それぞれの適用条件を扱う。

\section{分布モデルとパラメータ推定}

\subsection{NumPy 配列を入力とする専門計算機能}

SciPy の多くの関数は NumPy 配列を入力とし、NumPy 配列またはスカラーを返す。したがって、入力の \code{shape} と \code{dtype}、返り値の意味を NumPy の配列操作と対応付ける。目的と subpackage の対応は次のとおりである。

\begin{itemize}
\item \code{scipy.stats}\index{scipy@SciPy!stats@\code{scipy.stats}|textbf}: 分布、検定、要約統計
\item \code{scipy.optimize}: 最適化、非線形最小二乗、\code{curve\_fit}\index{scipy@SciPy!curve\_fit@\code{curve\_fit}|textbf}
\item \code{scipy.interpolate}\index{scipy@SciPy!interpolate@\code{scipy.interpolate}|textbf}: 補間
\item \code{scipy.signal}\index{scipy@SciPy!signal@\code{scipy.signal}|textbf}: フィルタ、ピーク検出、相関、スペクトル処理
\item \code{scipy.integrate}\index{scipy@SciPy!integrate@\code{scipy.integrate}|textbf}: 数値積分
\item \code{scipy.constants}\index{scipy@SciPy!constants@\code{scipy.constants}|textbf}: 物理定数
\end{itemize}

要素ごとの演算や基本的な集計には NumPy を用い、確率分布、最適化、補間、信号処理などの専用アルゴリズムが必要な場合に、対応する SciPy の subpackage を選ぶ。

\subsection{関数形の仮定とパラメータ推定}

分布の幅、減衰曲線の時定数、波形のピーク時刻を推定するには、それぞれモデル推定、非線形最小二乗、ピーク検出の条件を定める必要がある。SciPy は、これらの計算と推定結果の評価に用いる関数を提供する。

フィットでは、仮定した関数形のパラメータをデータから推定する。関数形の選択には物理的または統計的な根拠が必要であり、推定値の解釈にはフィット範囲、ビン幅、外れ値処理、残差、不確かさの確認が必要となる。

フィットの目的は、仮定したモデルのパラメータをデータから推定することである。分布の中心、幅、時定数などを比較可能な量として得るには、推定値だけでなく、不確かさ、フィット範囲、残差、モデルの適用条件も確認する。

\subsection{正規分布と指数分布}

多数の独立な微小変動の和は、条件を満たせば正規分布で近似できる。一定の発生率を持つ Poisson 過程では、事象間の待ち時間が指数分布に従う。対象データがこれらの仮定を満たすかは、生成過程、残差、代替モデルとの比較によって検証する。

たとえば、平均 \(\mu\)、標準偏差 \(\sigma\) のガウシアンは
\[
p(x) = \frac{1}{\sqrt{2\pi}\sigma}
\exp\!\left[-\frac{(x-\mu)^2}{2\sigma^2}\right]
\]
で書ける。また、時定数 \(\tau\) の指数分布は
\[
p(t) = \frac{1}{\tau} \exp\!\left(-\frac{t}{\tau}\right)
\qquad (t \ge 0)
\]
である。SciPy の分布オブジェクトは、確率密度、累積分布、分位点、乱数生成を共通のインタフェースで提供する。

\section{統計・フィット・補間・信号処理}

\subsection{\code{scipy.stats} による要約統計と分布}

配列の要素数、中心、広がり、形状を数値で確認する方法を示す。

\begin{lstlisting}[language=Python]
import numpy as np
from scipy import stats

values = np.array([1.2, 0.9, 1.5, 1.1, 0.8])

summary = stats.describe(values)
print(summary.nobs)
print(f"{summary.mean:.3f}")
print(f"{summary.variance:.3f}")
\end{lstlisting}

この例では、要素数、平均、分散が順に表示される。

\begin{lstlisting}[numbers=none, xleftmargin=2\zw]
5
1.100
0.075
\end{lstlisting}

\code{stats.describe()} は、要素数、平均、分散、最小値、最大値、歪度、尖度を名前付きタプルとして返す。分散の既定は \code{ddof=1} であり、NumPy の \code{var()} の既定とは異なる。各統計量の定義と単位を確認してから利用する。

SciPy の分布オブジェクトは、確率密度関数や累積分布関数を同じ API で扱える。

\begin{lstlisting}[language=Python]
x = np.linspace(-3.0, 3.0, 7)
pdf = stats.norm.pdf(x, loc=0.0, scale=1.0)
cdf = stats.norm.cdf(x, loc=0.0, scale=1.0)

print(pdf)
print(cdf)
\end{lstlisting}

\code{pdf} と \code{cdf} を同じ点列で評価すると、次のような値が得られる。

\begin{lstlisting}[numbers=none, xleftmargin=2\zw]
[0.00443185 0.05399097 0.24197072 0.39894228 0.24197072 0.05399097
 0.00443185]
[0.0013499  0.02275013 0.15865525 0.5        0.84134475 0.97724987
 0.9986501 ]
\end{lstlisting}

\code{pdf} は確率密度関数、\code{cdf} は累積分布関数、\code{rvs} は乱数生成、\code{ppf} は累積確率から分位点を求める逆関数に対応する。正規分布では \code{loc} は平均、\code{scale} は標準偏差であるが、これらの引数の意味は分布ごとに確認する。

ここで \code{pdf} は確率密度 \(p(x)\)、\code{cdf} はその積分
\[
F(x) = \int_{-\infty}^{x} p(t)\,dt
\]
に対応する。したがって、\code{stats.norm.cdf(x)} は「\(x\) 以下に入る確率」を返している。

\subsection{\code{curve\_fit} によるパラメータ推定}

\code{scipy.optimize.curve\_fit} は、利用者が定義したモデル関数のパラメータを非線形最小二乗法で推定する。\code{scipy.stats} が実装済みの確率分布に対する密度評価やパラメータ推定を提供するのに対し、\code{curve\_fit} には任意のモデル関数を渡せる。

\begin{lstlisting}[language=Python]
from scipy.optimize import curve_fit

def linear_model(x, a, b):
    return a * x + b

xdata = np.array([0.0, 1.0, 2.0, 3.0, 4.0])
ydata = np.array([1.1, 2.0, 3.2, 4.1, 5.1])

linear_params, linear_cov = curve_fit(linear_model, xdata, ydata)
linear_errors = np.sqrt(np.diag(linear_cov))

print(linear_params)
print(linear_errors)
\end{lstlisting}

この線形フィットでは、傾きと切片、および共分散行列から求めた標準誤差が次のように得られる。

\begin{lstlisting}[numbers=none, xleftmargin=2\zw]
[1.01 1.08]
[0.02516611 0.06164414]
\end{lstlisting}

\code{linear\_params} は推定パラメータ、\code{linear\_cov} は共分散行列である。対角成分の平方根が各パラメータの標準誤差に対応する。上の例では \code{linear\_params[0]} が傾き、\code{linear\_params[1]} が切片である。非対角成分はパラメータ間の共分散を表すため、相関を調べる場合は各パラメータの標準偏差で規格化する。\code{curve\_fit()} の出力は、仮定したモデル、残差の定義、\code{sigma} の指定に依存する。

上の直線モデル \(f(x; a, b) = ax + b\) なら、\code{curve\_fit()} は概ね
\[
S(a, b) = \sum_i \left[y_i - (a x_i + b)\right]^2
\]
を小さくする \(a\) と \(b\) を探していると考えればよい。より一般には、モデル \(f(x; \mathbf{p})\) に対して
\[
S(\mathbf{p}) = \sum_i \left[y_i - f(x_i; \mathbf{p})\right]^2
\]
という残差二乗和を最小化している。

得られた推定値は、データとモデルを重ねた図、別条件のデータ、単位、残差と照合する。SciPy の各関数についても、入力、返り値、仮定、失敗条件を確認する。

\section{モデル選択・重み付け・残差評価}

モデル推定では、次の項目を順に記録する。

\begin{enumerate}
\item 作図に基づく候補関数の検討
\item フィット範囲の決定
\item 必要に応じた妥当な初期値の設定
\item データとフィット曲線の重ね描き
\item 推定値と残差の確認
\end{enumerate}

モデルの定義とデータの確認を省いて \code{curve\_fit()} を実行しても、推定値の妥当性を判断できない。フィットは、モデルという仮説をデータと残差によって検証する手続きである。

\subsection{初期値、制約、残差}

モデルによっては、初期値や制約を与えないと不安定になる。

\begin{lstlisting}[language=Python]
def exp_model(x, a, tau, c):
    return a * np.exp(-x / tau) + c

x_decay = np.linspace(0.0, 8.0, 40)
rng = np.random.default_rng(31)
y_decay = 4.0 * np.exp(-x_decay / 2.5) + 0.3
y_decay += rng.normal(loc=0.0, scale=0.05, size=x_decay.size)

exp_params, exp_cov = curve_fit(
    exp_model,
    x_decay,
    y_decay,
    p0=[5.0, 2.0, 0.0],
    bounds=([0.0, 0.0, -np.inf], [np.inf, np.inf, np.inf]),
)
\end{lstlisting}

\code{p0}\index{p0@\code{p0}|textbf} は初期値、\code{bounds}\index{bounds@\code{bounds}|textbf} はパラメータ範囲である。どちらも収束を保証するものではないが、物理的にあり得ない範囲を除外する際に有効である。モデルによっては、初期値によって収束先が変わる。また、物理的に非負であるパラメータには、0 を下限とする制約を指定できる。

フィット結果の図には、関数形、推定パラメータ、データ系列を識別できる情報を記載する。凡例を用いる場合は、データや曲線を隠さない位置を選び、複数の図で同じ書式を用いる。

\begin{lstlisting}[language=Python]
import numpy as np
import matplotlib.pyplot as plt
from scipy import stats

rng = np.random.default_rng(37)
values = rng.normal(loc=3.0, scale=0.8, size=3000)
mean, sigma = stats.norm.fit(values)
x = np.linspace(values.min(), values.max(), 300)

fig, ax = plt.subplots()
ax.hist(values, bins=40, density=True, alpha=0.7, label="sample")
ax.plot(
    x,
    stats.norm.pdf(x, loc=mean, scale=sigma),
    linewidth=2.0,
    label=(
        "fit: "
        r"$f(x)=\frac{1}{\sqrt{2\pi}\sigma}"
        r"\exp\!\left[-\frac{(x-\mu)^2}{2\sigma^2}\right]$"
        "\n"
        + fr"$\mu={mean:.2f},\ \sigma={sigma:.2f}$"
    ),
)
ax.legend(loc="upper right")
\end{lstlisting}

図~\ref{fig:distribution-examples} では、フィットした分布名と推定パラメータを凡例へ記載している。左パネルは上の正規乱数と同じ入力を用い、右パネルは同じ乱数生成器で続けて生成した時定数 2.5 s、3000 個の指数乱数を用いる。本文には推定方法と解釈、キャプションには図の対象と条件、凡例には系列の識別情報を置く。

\begin{lstlisting}[language=Python]
model_y = linear_model(xdata, *linear_params)
residual = ydata - model_y
print(residual)
\end{lstlisting}

最適化が終了して推定値を返したことは、モデルの妥当性を保証しない。残差を横軸に対して表示し、符号や大きさに系統的な構造が残っていないか、特定の範囲だけ偏差が大きくないかを確認する。

外れ値やモデルの適用範囲外にある領域を含めると、推定値が変化する場合がある。フィット範囲と除外条件を推定値とともに記録する。

\begin{lstlisting}[language=Python]
mask = (xdata >= 0.0) & (xdata <= 3.0)
range_params, range_cov = curve_fit(
    linear_model,
    xdata[mask],
    ydata[mask],
)
\end{lstlisting}

フィット範囲は、モデルを適用する領域を定義する解析条件である。範囲をデータ確認後に変更した場合は、その判断基準と変更前後の結果も報告する。

\subsection{\code{sigma} を使った重み付きフィット}

物理の実データでは、各データ点の不確かさが等しいとは限らない。点ごとの標準偏差が既知であり、その大きさが異なる場合は、\code{sigma}\index{sigma@\code{sigma}|textbf} を使って残差を標準偏差で規格化する。\code{sigma} を与えないフィットはすべての点を等しい重みで扱うため、不確かさが点ごとに異なるデータでは目的と一致しない場合がある。

\begin{lstlisting}[language=Python]
y_err = np.array([0.2, 0.2, 1.5, 0.2, 0.2])

weighted_params, weighted_cov = curve_fit(
    linear_model,
    xdata,
    ydata,
    sigma=y_err,
    absolute_sigma=True,
)
\end{lstlisting}

この定義では、\code{sigma} が大きい点ほど規格化残差が小さくなり、目的関数への寄与が小さくなる。図~\ref{fig:fit-sigma} は、その差を模式的に示している。

\begin{figure}[htbp]
\centering
\includegraphics[width=0.7\linewidth]{fit_sigma_comparison.pdf}
\caption{\code{sigma} の有無によるフィット結果の違い。不確かさが点ごとに異なる場合、\code{sigma} の指定によって各残差の寄与が変わる。}
\label{fig:fit-sigma}
\end{figure}

\code{sigma} には、通常 \code{ydata} の各点に対応する 1 次元配列を渡す。このとき最小化される量は
\[
\chi^2(\mathbf{p}) = \sum_i
\left(
\frac{y_i - f(x_i; \mathbf{p})}{\sigma_i}
\right)^2
\]
と書けるので、\(\sigma_i\) が大きい点ほど目的関数への寄与は小さく、\(\sigma_i\) が小さい点ほど大きくなる。独立な Poisson 計数で平均値を観測数 \(N\) によって推定し、正規近似が妥当な範囲では、標準偏差を \(\sqrt{N}\) で近似する場合がある。低計数や重み付き計数には、その近似を機械的に適用しない。

\code{absolute\_sigma=True} は、共分散の計算に \code{sigma} の絶対的な大きさを使う指定である。渡した値が適切な標準偏差かどうかを、この設定が検証するわけではない。これを \code{False} のままにすると、\code{sigma} は相対的な重みとしては使われるが、共分散行列の大きさは残差のばらつきに合わせてあとから再スケーリングされる。そのため、推定値の重み付けだけでなく誤差評価まで行いたいなら、\code{sigma} の意味を明確にした上で \code{absolute\_sigma} の設定も意識する必要がある。

図~\ref{fig:fit-sigma} は、点ごとの不確かさを重みに反映した場合と、全点を等しい重みで扱った場合の違いを示している。どちらを採用するかは、\code{sigma} が測定の標準偏差を表すか、各点を同じ重みで比較する解析なのかによって決まる。

\code{curve\_fit()} の共分散とその対角成分から求める標準誤差は、解の近傍での線形近似に依存する~\cite{scipy_curvefit}。強い非線形性やパラメータ間の縮退がある場合には信頼できないことがある。1 次元の \code{sigma} による上の式は各点の誤差を独立とする。相関を扱うには共分散行列を渡す方法があり、通常の呼び出しは横軸 \code{xdata} の測定誤差を考慮しない。

\subsection{\code{scipy.interpolate} による補間}

補間は、測定点の間における値を仮定した関数形から推定する操作である。ここでは、1 次元データに対する区分線形補間を扱う。

\begin{lstlisting}[language=Python]
from scipy.interpolate import make_interp_spline

x = np.array([0.0, 1.0, 2.0, 4.0])
y = np.array([0.0, 2.0, 1.0, 3.0])

linear_interp = make_interp_spline(x, y, k=1)
x_new = np.array([0.5, 1.5, 3.0])
y_new = linear_interp(x_new, extrapolate=False)

print(y_new)
\end{lstlisting}

補間点 \code{x\_new} に対しては、次の値が返る。

\begin{lstlisting}[numbers=none, xleftmargin=2\zw]
[1.  1.5 2. ]
\end{lstlisting}

入力の \code{x} は重複がなく、昇順に並んでいる必要がある。区分線形補間は測定点を通る連続関数だが、接続点で傾きが不連続になり得る。上の \code{extrapolate=False} は入力範囲の外で \code{NaN} を返し、外挿を行わない指定である。範囲外まで同じ傾向が続くとは限らない。

線形補間では、各区間 \(x_i \le x \le x_{i+1}\) で
\[
y(x) = y_i + \frac{y_{i+1}-y_i}{x_{i+1}-x_i}(x-x_i)
\]
という一次式を使っている。したがって、測定点の間では直線でつながるが、測定範囲の外ではこの式をそのまま信じる根拠はない。

\subsection{\code{scipy.signal} による平滑化とピーク検出}

波形や時系列では、ピークの検出や軽い平滑化がよく必要になる。

\begin{lstlisting}[language=Python]
from scipy.signal import savgol_filter, find_peaks

x = np.linspace(0.0, 20.0, 200)
signal = np.sin(x) + 0.15 * np.random.default_rng(7).normal(size=x.size)

smooth = savgol_filter(signal, window_length=11, polyorder=3)
peaks, props = find_peaks(smooth, height=0.7, distance=15)

print(peaks)
print(np.round(props["peak_heights"], 3))
\end{lstlisting}

この例では、検出されたピーク位置とその高さが次のように表示される。

\begin{lstlisting}[numbers=none, xleftmargin=2\zw]
[ 13  76 138]
[0.977 1.048 0.994]
\end{lstlisting}

\code{savgol\_filter()} は局所多項式によって平滑化する。\code{find\_peaks()} は高さ、試料点間の距離、prominence などの条件でピーク候補を抽出する。検出条件をサンプリング間隔と対応付け、既知の合成波形で検出率と誤検出を確認してから実データへ適用する。

\subsection{\code{scipy.integrate} と \code{scipy.constants}}

\code{scipy.integrate} は離散点または関数の数値積分を、\code{scipy.constants} は物理定数の数値と単位情報を提供する。

\begin{lstlisting}[language=Python]
from scipy import constants
from scipy.integrate import simpson

x = np.linspace(0.0, np.pi, 101)
y = np.sin(x)

area = simpson(y, x=x)
print(area)
print(constants.c)
\end{lstlisting}

この積分値と光速定数は、次のように出力される。

\begin{lstlisting}[numbers=none, xleftmargin=2\zw]
2.0000000108245044
299792458.0
\end{lstlisting}

\code{simpson()} は、ここでは
\[
\int_0^\pi \sin x\,dx = 2
\]
という積分を、離散点 \((x_i, y_i)\) から数値的に評価している。したがって、\code{simpson(y, x=x)} は「\(y(x)\) をサンプル点から積分する関数」と読めばよい。また、\code{constants.c} は
\[
c = 299792458\ \mathrm{m\,s^{-1}}
\]
を表す定数である。\code{constants.c} 自体は単位を保持しない浮動小数点数なので、単位を演算中も保持する必要がある場合は、次章の Astropy の \code{Quantity} と定数オブジェクトを用いる。

\section{モデル・制約・外挿に関する注意点}

\subsection{モデル選択の理由がないままフィットしない}

ガウシアン、指数、直線の選択には、物理的または統計的な仮定が必要である。曲線がデータ点の近くを通ることだけでは、モデルの生成過程やパラメータの解釈を正当化できない。

\subsection{初期値と制約を無視しない}

非線形フィットは、初期値次第で別の極小へ落ちたり、収束しなかったりする。\code{curve\_fit()} が値を返したからといって、それが唯一の解だとは限らない。複数の初期値で試す、図へ重ねて見る、残差を見る、といった確認が必要である。

\subsection{補間と外挿を混同しない}

\code{make\_interp\_spline()} などの補間関数は、観測点の間を埋めるための道具である。測定範囲の外まで同じ関係が続く根拠は補間からは得られない。外側を予測する場合は、外挿に用いるモデルと適用範囲を別に検討する。

\subsection{信号処理におけるサンプリング間隔}

\code{find\_peaks()} の \code{distance} やフィルタ長は、サンプル数で指定することが多い。したがって、時間軸が 1 サンプル 1 ns なのか 10 ns なのかで、同じ数値でも意味が変わる。コード中の数字をそのまま流用せず、軸の単位へ戻して考える必要がある。

\section{発生間隔分布のモデル比較}

SciPy は、解析の中で次のような役割を持つ。

\begin{enumerate}
\item \code{stats} による分布形状と要約統計の確認
\item \code{signal} によるピーク候補の抽出と平滑化
\item \code{interpolate} による校正曲線・補助曲線の作成
\item \code{optimize} によるモデルパラメータの推定
\item \code{integrate} と \code{constants} による計算の完結
\end{enumerate}

SciPy は、入力配列の物理的意味や、採用すべきモデルを自動では決めない。特にフィットでは、データ範囲、重み、残差、推定値の不確かさを確認し、採用した条件を結果とともに記録する必要がある。

ここまで見てきた SciPy は、主に「数値をどう計算するか」を受け持つ道具である。時刻系、単位、座標系、FITS のように「数値そのものの意味付け」まで含めて扱いたくなったところで、次章の Astropy が必要になる。SciPy は計算アルゴリズム、Astropy は物理量や観測量の表現を担当すると考えると、章のつながりも見えやすい。

本文を読み終えた後で関数名や引数を確認したいときは、SciPy User Guide~\cite{scipy_user_guide} を起点にすると探しやすい。分布や検定は Statistics~\cite{scipy_stats}、フィットや最適化は Optimization~\cite{scipy_optimize}、補間は Interpolation~\cite{scipy_interpolate}、波形処理は Signal processing~\cite{scipy_signal}、数値積分は Integration~\cite{scipy_integrate} が対応している。本章は、それらを解析で使う順に並べ替えて説明した。

\section{章末課題：分布モデルとフィットの演習}

この章の課題は、既知の分布による API と結果の読み方の検証、および実データに対するモデルの検討という二段階からなる。次のコードが検証用データを生成する。乱数シードを固定しているため、同じライブラリ環境では同じ入力を再現できる。

\begin{lstlisting}[language=Python]
import numpy as np

rng = np.random.default_rng(20260412)
synthetic_dt = rng.exponential(scale=120.0, size=5000)
x = np.linspace(0.0, 10.0, 80)
y_true = 2.5 * x + 1.2
sigma = np.full_like(x, 0.5)
y_obs = y_true + rng.normal(scale=sigma)
\end{lstlisting}

実データに当てる分布と関数を、最初から決め打ちしてはならない。まず図を描いて特徴を観察し、比較する量と近似を試す範囲を判断せよ。ヒストグラムと累積分布の選択、短時間側の取扱い、対数軸の採否についても、その理由とともに記録せよ。

\begin{enumerate}
\item \textbf{既知の指数分布による推定方法の検証}。
\code{scipy.stats.expon.fit(synthetic\_dt, floc=0)} で指数分布のパラメータを推定せよ。\code{floc=0} は位置母数を 0 に固定する指定であり、返り値は \code{loc} と \code{scale} の順である。推定された \code{scale} が生成時の 120 秒に近いことを確認し、標本平均とも比較せよ。許容範囲を 115 秒から 125 秒として自動確認する \code{assert} も記述せよ。

\item \textbf{重み付き直線フィットと残差の確認}。
\code{curve\_fit()} を使い、\code{y\_obs} を \(y=ax+b\) へフィットせよ。\code{sigma=sigma}、\code{absolute\_sigma=True} を指定し、推定値 \code{popt} と共分散行列 \code{pcov} を受け取る。傾きが 2.3 から 2.7、切片が 0.7 から 1.7 の範囲に入ることを確認せよ。\code{np.sqrt(np.diag(pcov))} が各パラメータの標準不確かさを与えること、\code{y\_obs - model(x, *popt)} が残差になることを説明し、データ、フィット線、残差を 2 段の図へ保存せよ。

\item \textbf{実データの発生間隔分布と fit 方針の決定}。
第~\ref{chap:numpy}~章で作った \(\Delta t\) 配列がある場合は、等間隔ビンのヒストグラム、対数ビンのヒストグラム、経験累積分布のうち少なくとも 2 通りを作成せよ。その上で、指数分布を含む候補を比較し、採用するモデルとフィット範囲を決定せよ。1 問目の検証用データと異なり、実データが単一の指数分布に従う保証はない。モデルを採用した根拠、除外した範囲、残差または適合度を記録せよ。

\item \textbf{解析上の選択に対する感度評価}。
3 問目と少なくとも 1 条件を変えて推定を繰り返せ。ビン幅、フィット範囲、短時間側の扱い、重み、関数形などから選び、推定値と不確かさがどの程度変化したかを表にまとめよ。結果が変わったという事実だけでなく、どのデータ点の寄与が変化したためかを図と残差から説明せよ。

実データでは、検索期間・地域、使用したイベント種別、マグニチュードの型を記録する。USGS の CSV には \code{magType} などの列があり~\cite{usgs_csv}、異なる尺度や検出能力の地域を混ぜると単一のモデルで説明できない場合がある。

\item \textbf{発展課題：マグニチュード分布と Gutenberg--Richter 関係}。
USGS カタログのマグニチュード \(M\) について、度数分布と累積件数 \(N(\geq M)\) を作り、カタログが十分に完全とみなせる下限 \(M_c\) を検討せよ。\(M\geq M_c\) の範囲で Gutenberg--Richter 関係~\cite{gutenberg_richter_1944}
\[
\log_{10}N(\geq M)=a-bM
\]
が成り立つかを調べる。この関係は、地震規模に対するべき乗則を、対数尺度であるマグニチュード \(M\) で表したものである。

累積件数の点は同じイベントを共有するため、独立な Poisson 変数として \(\chi^2\) を解釈できない。推定には、重ならないマグニチュード区間の度数に対する計数モデル、または個々の \(M\) に対する尤度法を用いよ。累積分布は表示だけに用い、推定量と区別せよ。日本近海または全球を対象とし、期間、空間範囲、\(M_c\)、ビン、イベント数、推定法、不確かさ、残差を明示せよ。
\end{enumerate}

SciPy は数値配列に対する計算を担うが、物理単位、時刻系、天球座標の意味までは自動的に補わない。次章では、Astropy によって数値へ単位や座標系を保持し、SciPy や表処理ライブラリとの間で受け渡す方法を扱う。
